% !TEX program = xelatex
% AI-effects 프로젝트 --- AI 점유율 지표 A(대화 점유율)와 B(절감 시간 점유율)의 순위 격차가 큰 직업.
% 표 본문: results/table/31_summary.tex, 31_top_pos.tex, 31_top_neg.tex, 31_top30A.tex, 31_ksco1.tex
%          (code/04_14_analysis_kor_conv_saved_gap.do가 생성)
% 2026-10-24 재설계: 지표 A·B를 AEI 업무용 대화 비중에서 직접 만듦(총량·K 없음, 두 달 = 월별 비중 단순평균). 수치 갱신.
% 전체 자료: results/table/31_kor_conv_saved_gap.xlsx
% 관련: paper 25(상대 활용도, tau_eff), paper 26(소분류 × SOC 구성), paper 27(점유율 지표 A~C), paper 29(점유율과 노출도의 격차)
% 컴파일: xelatex -> biber -> xelatex -> xelatex (kotex, biblatex/biber 필요)
\documentclass[11pt]{article}

\usepackage{fontspec}
\usepackage{kotex}
\usepackage[a4paper,margin=2.0cm]{geometry}
\usepackage{longtable}
\usepackage{booktabs}
\usepackage{array}
\usepackage{amsmath}
\usepackage{xcolor}
\usepackage{hyperref}
\usepackage{enumitem}
\usepackage[backend=biber,style=authoryear,maxcitenames=2,uniquelist=false,sorting=nyt]{biblatex}
\addbibresource{../../reference/references.bib}

\hypersetup{colorlinks=true,linkcolor=blue!50!black,urlcolor=blue!50!black}
\newcolumntype{L}[1]{>{\raggedright\arraybackslash}p{#1}}
\newcolumntype{R}[1]{>{\raggedleft\arraybackslash}p{#1}}
\setlength{\LTleft}{0pt}
\setlength{\LTright}{0pt}
\renewcommand{\arraystretch}{1.15}
\setlength{\tabcolsep}{3.5pt}

\title{대화 점유율 지표(A)와 절감 시간 지표(B)의 순위 격차가 큰 직업\\[2mm]
\large --- 대화당 절감 시간과 SOC 구성으로 본 차이 ---}
\author{AI-effects 프로젝트}
\date{\today}

\begin{document}
\maketitle

\begin{abstract}
\noindent paper 27의 AI 점유율 지표 중 A(업무용 대화 점유율)와 B(AI 절감 시간 점유율)의 순위가 크게 다른 직업을 KSCO 8차 소분류 129개에서 찾았다. 격차는 순위 격차 $\mathrm{RG} = r(\text{A}) - r(\text{B})$로 쟀다. 양수이면 절감 시간 순위가 대화 점유율 순위보다 높은 직업이다. B/A는 그 직업의 대화당 절감 시간 $\tau$에 비례하므로, 두 지표의 차이는 모두 $\tau$에서 나온다. 그리고 SOC별 대화당 절감 시간은 AEI의 전 세계(GLOBAL) 값이므로, 한국 소분류 사이의 $\tau$ 차이는 그 소분류로 들어온 대화의 SOC 구성에서 나온다.
\begin{itemize}[leftmargin=*]
  \item \textbf{두 지표는 대체로 같다.} A와 B의 Spearman 순위상관은 0.96이다. 대화 점유율 상위 10개 중 8개가 절감 시간 상위 10개에도 있다. 순위 격차의 절댓값은 평균 7.4계단이고, 20계단 이상인 소분류는 9개다.
  \item \textbf{대화당 절감 시간 $\tau$는 소분류마다 크게 다르다.} 전체 평균은 259분이고, 소분류 값은 21$\sim$473분(90분위/10분위 6.9배)이다. $\tau$는 노출도와 양의 상관(0.36$\sim$0.47), 취업자 수와 음의 상관($-0.36$)을 보인다.
  \item \textbf{B 순위가 더 높은 직업(RG > 0)}: 226 통신 및 방송 송출 장비 기사(+21), 159 기타 판매 및 고객 서비스 관리자(+18), 132 보험 및 금융 관리자, 139 기타 전문 서비스 관리자, 951 가사 및 육아 도우미(각 +16). 상위 15개 중 7개가 관리자(대분류 1)이고 4개가 돌봄$\cdot$교육 직업이다. 대화가 대화당 400$\sim$540분짜리 관리자 SOC나 돌봄$\cdot$교육 SOC로 연결된 직업이다.
  \item \textbf{A 순위가 더 높은 직업(RG < 0)}: 411 경찰$\cdot$소방 및 교도 관련 종사자($-32$), 992 계기$\cdot$자판기 및 주차 관리 종사자, 999 기타 서비스 관련 단순 종사자(각 $-30$), 522 상품 대여 종사자($-26$), 521 매장 판매 종사자($-25$). 대화가 대화당 20$\sim$90분짜리 판매$\cdot$서비스$\cdot$보안 SOC로 연결된 직업이다. 대화 점유율 상위 30에 드는 522와 521은 B로 보면 각각 48위, 55위로 내려간다.
  \item \textbf{대분류}: 관리자는 B 점유율이 A보다 2.0\%p 높고(16개 중 13개 RG > 0), 판매 종사자는 1.9\%p 낮다(5개 모두 RG < 0). 판매$\cdot$서비스$\cdot$단순노무$\cdot$농림어업은 $\tau$가 33$\sim$122분으로 전체 평균의 절반 이하다.
\end{itemize}
\end{abstract}

\section{두 지표가 갈리는 이유}

paper 27의 정의를 그대로 쓴다. 소분류 $j$의 월평균 업무용 대화 수 $C_j$와 AI 절감 시간 $S_j$(시간)에 대해
\[
\text{A}_j = 100 \times \frac{C_j}{\sum_k C_k}, \qquad \text{B}_j = 100 \times \frac{S_j}{\sum_k S_k}
\]
다. 소분류의 대화당 실효 절감 시간을 $\tau_j = 60 \times S_j / C_j$(분, paper 25의 \texttt{tau\_eff})라 하고 전체 값을 $\bar\tau = 60 \times \sum_k S_k / \sum_k C_k$라 하면
\[
\frac{\text{B}_j}{\text{A}_j} = \frac{\tau_j}{\bar\tau}, \qquad \ln \text{B}_j = \ln \text{A}_j + \ln(\tau_j / \bar\tau)
\]
다. 즉 B는 A에 ``대화당 절감 시간이 평균보다 얼마나 긴가''를 곱한 값이다. 두 지표의 순위가 갈리는 이유는 $\tau_j$ 하나뿐이다. 이 129개 소분류에서는 모든 대화에 절감 시간 값이 있으므로(paper 25의 \texttt{ts\_cover} = 1) 절감 시간이 없어 생기는 차이는 없다.

$\tau_j$는 다시 SOC 단계로 나뉜다. $S_j$는 SOC $s$의 한국 대화 수에 그 SOC의 대화당 절감 시간 $\tau_s$를 곱한 뒤 연계표 몫(paper 21)으로 소분류에 배분한 값이다(paper 25$\cdot$26과 같은 계산). 그래서
\[
\tau_j = \sum_s \frac{c_{sj}}{C_j} \, \tau_s
\]
다. $c_{sj}$는 SOC $s$가 소분류 $j$에 주는 대화 수다. $\tau_s$는 AEI release\_2026\_06\_26의 세부 직업별 \texttt{human\_only\_time} $-$ \texttt{human\_with\_ai\_time}이다. 이 값은 Claude가 대화마다 추정한 시간의 평균이다\parencite{tamkin-2025-estimating-ai-productivity-gains-from,appel-2026-economic-primitives}. 세부 직업의 시간 지표는 GLOBAL에만 공개되므로, 한국 대화에도 전 세계 값을 쓴다. 결국 소분류 사이 $\tau_j$의 차이는 \emph{그 소분류로 들어온 한국 대화가 어떤 SOC의 대화인가}에서만 생긴다.

어느 SOC가 $\tau_j$를 올리거나 내렸는지는 SOC별 기여
\[
d_{sj} = 100 \times \left( \frac{v_{sj}}{S_j} - \frac{c_{sj}}{C_j} \right) \quad (\%\text{p})
\]
로 본다. $v_{sj}$는 SOC $s$가 소분류 $j$에 주는 절감 시간이다. 절감 시간 몫이 대화 몫보다 큰 SOC($d_{sj} > 0$)가 $\tau_j$를 올린다. $\sum_s d_{sj} = 0$이고, SOC가 하나뿐인 소분류는 $d = 0$이다.

\section{격차 지표}

\begin{itemize}[leftmargin=*]
  \item \textbf{순위 격차} $\mathrm{RG}_j = r(\text{A}_j) - r(\text{B}_j)$. $r(\cdot)$은 129개 안 순위(1 = 최상위, 같은 값은 소분류 코드 순)다. 양수이면 B 순위가 높다. 이 노트의 주 지표다.
  \item \textbf{로그 비율} $\mathrm{LR}_j = \ln(\text{B}_j / \text{A}_j) = \ln(\tau_j / \bar\tau)$. 순위 격차와 달리 직업 규모와 순위 위치에 영향받지 않는 강도 비교다. $\tau_j$ 순위와 같다.
  \item \textbf{점유율 차이} $\mathrm{DS}_j = \text{B}_j - \text{A}_j$(\%p, 합 0). 큰 직업일수록 커진다.
\end{itemize}
순위 격차는 같은 $\tau_j$라도 분포 위치에 따라 크기가 달라진다. A 분포의 위쪽은 값의 간격이 넓어 순위가 잘 움직이지 않고, 가운데는 간격이 좁아 조금만 달라도 여러 계단 움직인다. 실제로 $|\mathrm{RG}|$ 평균은 A 순위 1$\sim$43위 5.3, 44$\sim$86위 11.7, 87$\sim$129위 5.3이다(표~\ref{tab:summary}). 아래쪽(87위 이하)은 순위가 129위에서 막혀 아래로 덜 움직인다. 그래서 RG와 LR의 순위상관은 0.78에 그친다. 이 노트의 목록은 RG로 만들되, 표에 $\tau_j$를 함께 적어 강도 차이를 바로 볼 수 있게 했다.

\section{결과}

\subsection{요약}

\begin{longtable}{L{12.5cm}R{2.5cm}}
\caption{A와 B의 관계, 대화당 절감 시간, 순위 격차 요약(소분류 129개)}\label{tab:summary}\\
\toprule
항목 & 값 \\
\midrule
\endfirsthead
% 04_14_analysis_kor_conv_saved_gap.do가 생성. 직접 고치지 말 것.
A와 B의 Spearman 순위상관 & 0.96 \\
\(\ln\) A와 \(\ln\) B의 Pearson 상관 & 0.97 \\
전체 대화당 절감 시간 \(\bar\tau\)(분) & 259.4 \\
소분류 \(\tau\) 중앙값(분) & 196.8 \\
소분류 \(\tau\) 최솟값(분) & 21.2 \\
소분류 \(\tau\) 최댓값(분) & 473.4 \\
소분류 \(\tau\) 90분위 / 10분위(배) & 6.89 \\
\(|\mathrm{RG}|\) 최댓값 & 32 \\
\(|\mathrm{RG}|\) 평균 & 7.4 \\
\(|\mathrm{RG}| \ge 20\)인 소분류 수 & 9 \\
\(|\mathrm{RG}|\) 평균: A 순위 1\(\sim\)43위 & 5.3 \\
\(|\mathrm{RG}|\) 평균: A 순위 44\(\sim\)86위 & 11.7 \\
\(|\mathrm{RG}|\) 평균: A 순위 87\(\sim\)129위 & 5.3 \\
RG와 LR의 Spearman 순위상관 & 0.78 \\
RG와 취업자 수의 Spearman 순위상관 & -0.38 \\
\(\tau\)와 취업자 수의 Spearman 순위상관 & -0.36 \\
\(\tau\)와 Observed exposure의 Spearman 순위상관 & 0.36 \\
\(\tau\)와 Felten LM AIOE의 Spearman 순위상관 & 0.47 \\
\(\tau\)와 Eloundou GPT-4 \(\beta\)의 Spearman 순위상관 & 0.46 \\

\bottomrule
\end{longtable}

A와 B는 순위상관 0.96으로 거의 같은 순서를 준다. 그러나 $\tau_j$는 소분류 사이에 6.9배(90분위/10분위)나 차이가 난다. 그런데도 순위상관이 높은 이유는 A 자체의 분포가 훨씬 넓기 때문이다(A는 0.0007\%$\sim$11.9\%, 약 16,000배). 소분류 $\tau_j$의 단순 평균은 217분, 중앙값은 197분으로, 대화 점유율 가중 평균 $\bar\tau$(259분)보다 짧다. 대화가 많은 직업일수록 대화당 절감 시간이 긴 편이라는 뜻이다. $\ln \text{B} = \ln \text{A} + \ln(\tau/\bar\tau)$에서 $\ln \text{A}$의 분산이 $\ln \tau$의 분산보다 훨씬 커서 순서가 크게 바뀌지 않는다.

순위 격차가 양수인 소분류는 68개, 음수는 52개, 0은 9개다. 양수 쪽은 작은 격차가 많고, 음수 쪽은 크기가 크다($|\mathrm{RG}| \ge 20$ 9개 중 8개가 음수). 순위 격차는 로그 비율 $\mathrm{LR}_j = \ln(\tau_j/\bar\tau)$과 같은 방향으로 움직이는데, $\tau_j/\bar\tau$는 0.08$\sim$1.83배라 LR은 $-2.50\sim+0.60$이다. 즉 로그 척도에서 아래쪽 꼬리가 훨씬 길다. 이는 ``대화는 많지만 대화당 절감 시간이 아주 짧은'' 직업이 소수 있고, 이들이 B에서 크게 밀린다.

$\tau_j$는 노출도 지표와 양의 상관(Observed exposure 0.36, Felten 0.47, Eloundou $\beta$ 0.46)을 보이고, 취업자 수와는 음의 상관($-0.36$)을 보인다. 노출도가 높은 사무$\cdot$전문 직업의 대화가 대화당 절감 시간도 길다는 뜻이다. 다만 그 효과는 작다. paper 28에서 노출도와의 순위상관은 B가 A보다 Felten(0.63 대 0.60)과 Eloundou $\beta$(0.70 대 0.68)에서 약간 높고, Observed exposure(0.69 대 0.70)에서는 약간 낮다.

\subsection{순위 격차가 큰 소분류}

표~\ref{tab:pos}와 표~\ref{tab:neg}의 마지막 열은 그 소분류의 $\tau_j$를 가장 많이 올린(표~\ref{tab:pos}) 또는 내린(표~\ref{tab:neg}) SOC, 그 SOC의 대화당 절감 시간 $\tau_s$, 기여 $d_{sj}$다.

\begin{small}
\begin{longtable}{R{0.5cm}L{3.6cm}R{0.9cm}R{0.7cm}R{0.9cm}R{0.7cm}R{0.8cm}R{0.8cm}L{6.0cm}}
\caption{순위 격차 상위 15: B(절감 시간) 순위가 A(대화 점유율) 순위보다 높은 소분류}\label{tab:pos}\\
\toprule
 & 소분류 & A(\%) & $r$(A) & B(\%) & $r$(B) & RG & $\tau$(분) & $\tau$를 가장 올린 SOC($\tau_s$, $d_{sj}$) \\
\midrule
\endfirsthead
\toprule
 & 소분류 & A(\%) & $r$(A) & B(\%) & $r$(B) & RG & $\tau$(분) & $\tau$를 가장 올린 SOC($\tau_s$, $d_{sj}$) \\
\midrule
\endhead
% 04_14_analysis_kor_conv_saved_gap.do가 생성. 직접 고치지 말 것.
1 & 226 통신 및 방송 송출 장비 기사 & 0.10 & 84 & 0.17 & 63 & +21 & 434 & 27-4011 Audio and Video Technicians (450분, +3.1) \\
2 & 159 기타 판매 및 고객 서비스 관리자 & 0.09 & 86 & 0.13 & 68 & +18 & 383 & 11-1021 General and Operations Managers (541분, +6.9) \\
3 & 132 보험 및 금융 관리자 & 0.22 & 61 & 0.40 & 45 & +16 & 473 & 11-1021 General and Operations Managers (541분, +1.9) \\
4 & 139 기타 전문 서비스 관리자 & 0.24 & 59 & 0.41 & 43 & +16 & 434 & 11-2021 Marketing Managers (499분, +10.8) \\
5 & 951 가사 및 육아 도우미 & 0.11 & 83 & 0.13 & 67 & +16 & 306 & 39-9011 Childcare Workers (334분, +1.8) \\
6 & 151 판매 및 운송 관리자 & 0.34 & 51 & 0.59 & 36 & +15 & 451 & 11-2021 Marketing Managers (499분, +5.5) \\
7 & 941 청소 관련 종사자 & 0.04 & 99 & 0.05 & 84 & +15 & 299 & 37-1011 First-Line Supervisors of Housekeeping and Janitorial Workers (303분, +1.5) \\
8 & 112 기업 대표 및 기업 고위 임원 & 0.15 & 72 & 0.23 & 58 & +14 & 396 & 11-1021 General and Operations Managers (541분, +6.9) \\
9 & 421 교사 보조 및 아동 돌봄 종사자 & 0.09 & 85 & 0.11 & 71 & +14 & 341 & 25-9042 Teaching Assistants, Preschool, Elementary, Middle, and Secondary School, Except Special Education (378분, +4.8) \\
10 & 851 금속 공작 기계 조작원 & 0.05 & 97 & 0.05 & 83 & +14 & 250 & 51-4041 Machinists (249분, +0.0) \\
11 & 251 보육교사 & 0.02 & 110 & 0.03 & 97 & +13 & 326 & 39-9011 Childcare Workers (334분, +2.3) \\
12 & 262 학교 교사 & 0.14 & 74 & 0.19 & 61 & +13 & 346 & 25-2051 Special Education Teachers, Preschool (471분, +13.8) \\
13 & 322 운송 및 무역 사무원 & 0.05 & 98 & 0.05 & 85 & +13 & 270 & 41-1012 First-Line Supervisors of Non-Retail Sales Workers (369분, +0.2) \\
14 & 141 건설‧전기 및 생산 관련 관리자 & 0.54 & 40 & 0.83 & 28 & +12 & 400 & 11-3051 Industrial Production Managers (413분, +2.3) \\
15 & 149 기타 건설‧전기 및 생산 관련 관리자 & 0.42 & 46 & 0.65 & 34 & +12 & 400 & 11-3051 Industrial Production Managers (413분, +3.0) \\

\bottomrule
\end{longtable}

\begin{longtable}{R{0.5cm}L{3.6cm}R{0.9cm}R{0.7cm}R{0.9cm}R{0.7cm}R{0.8cm}R{0.8cm}L{6.0cm}}
\caption{순위 격차 하위 15: A(대화 점유율) 순위가 B(절감 시간) 순위보다 높은 소분류}\label{tab:neg}\\
\toprule
 & 소분류 & A(\%) & $r$(A) & B(\%) & $r$(B) & RG & $\tau$(분) & $\tau$를 가장 내린 SOC($\tau_s$, $d_{sj}$) \\
\midrule
\endfirsthead
\toprule
 & 소분류 & A(\%) & $r$(A) & B(\%) & $r$(B) & RG & $\tau$(분) & $\tau$를 가장 내린 SOC($\tau_s$, $d_{sj}$) \\
\midrule
\endhead
% 04_14_analysis_kor_conv_saved_gap.do가 생성. 직접 고치지 말 것.
1 & 411 경찰‧소방 및 교도 관련 종사자 & 0.38 & 47 & 0.08 & 79 & -32 & 56 & 33-3012 Correctional Officers and Jailers (33분, -35.1) \\
2 & 992 계기‧자판기 및 주차 관리 종사자 & 0.13 & 76 & 0.01 & 106 & -30 & 24 & 39-3093 Locker Room, Coatroom, and Dressing Room Attendants (24분, +0.0) \\
3 & 999 기타 서비스 관련 단순 종사자 & 0.34 & 52 & 0.06 & 82 & -30 & 44 & 39-3093 Locker Room, Coatroom, and Dressing Room Attendants (24분, -35.5) \\
4 & 522 상품 대여 종사자 & 1.07 & 22 & 0.37 & 48 & -26 & 90 & 41-2021 Counter and Rental Clerks (90분, 0.0) \\
5 & 521 매장 판매 종사자 & 0.83 & 30 & 0.26 & 55 & -25 & 80 & 41-2011 Cashiers (58분, -22.4) \\
6 & 422 요양보호사 및 간병인 & 0.11 & 82 & 0.02 & 104 & -22 & 37 & 31-1131 Nursing Assistants (37분, +0.0) \\
7 & 532 방문 및 노점 판매 종사자 & 0.12 & 81 & 0.02 & 103 & -22 & 36 & 41-9091 Door-to-Door Sales Workers, News and Street Vendors, and Related Workers (36분, -0.0) \\
8 & 241 의사·한의사 및 치과의사 & 0.15 & 71 & 0.04 & 91 & -20 & 66 & 29-1212 Cardiologists (48분, -20.5) \\
9 & 135 정보 통신 관련 관리자 & 0.54 & 39 & 0.24 & 57 & -18 & 115 & 11-3021 Computer and Information Systems Managers (115분, 0.0) \\
10 & 311 기획 및 마케팅 사무원 & 0.21 & 63 & 0.06 & 81 & -18 & 81 & 43-4151 Order Clerks (53분, -6.9) \\
11 & 953 판매 관련 단순 종사자 & 0.13 & 79 & 0.03 & 96 & -17 & 63 & 53-7065 Stockers and Order Fillers (44분, -9.1) \\
12 & 331 회계 및 경리 사무원 & 0.44 & 45 & 0.23 & 59 & -14 & 135 & 43-3031 Bookkeeping, Accounting, and Auditing Clerks (125분, -6.4) \\
13 & 249 기타 보건 전문가 및 관련 종사자 & 0.30 & 53 & 0.14 & 66 & -13 & 118 & 29-1212 Cardiologists (48분, -7.6) \\
14 & 531 통신 관련 판매 종사자 & 0.22 & 60 & 0.11 & 72 & -12 & 128 & 41-9091 Door-to-Door Sales Workers, News and Street Vendors, and Related Workers (36분, -6.3) \\
15 & 246 치료‧재활사 및 의료기사 & 0.57 & 38 & 0.37 & 49 & -11 & 168 & 19-3033 Clinical and Counseling Psychologists (148분, -6.6) \\

\bottomrule
\end{longtable}
\end{small}

\paragraph{B 순위가 높은 직업.} 상위 15개의 $\tau_j$는 249$\sim$473분으로 모두 소분류 중앙값(197분)보다 길다. 세 묶음으로 나뉜다.
\begin{itemize}[leftmargin=*]
  \item \textbf{관리자}(159, 132, 139, 151, 112, 141, 149). 대화가 11-1021 General and Operations Managers(541분), 11-2021 Marketing Managers(499분), 11-3031 Financial Managers(463분), 11-3051 Industrial Production Managers(413분)처럼 대화당 절감 시간이 긴 관리자 SOC로 연결된다. AEI는 관리자 SOC의 대화를 사람이 혼자 하면 오래 걸리는 과업으로 추정한다는 뜻이다(어떤 과업인지는 이 노트에서 확인하지 않았다).
  \item \textbf{돌봄$\cdot$교육}(951 가사 및 육아 도우미, 421 교사 보조 및 아동 돌봄, 251 보육교사, 262 학교 교사). 39-9011 Childcare Workers(334분), 25-9042 Teaching Assistants(378분), 25-2051 Special Education Teachers, Preschool(471분)로 연결된다. 대화 수는 적어도 대화당 절감 시간이 길다.
  \item \textbf{그 밖의 직업}(226 통신 및 방송 송출 장비 기사, 851 금속 공작 기계 조작원, 941 청소 관련 종사자, 322 운송 및 무역 사무원). 226은 대화의 84\%가 27-4011 Audio and Video Technicians(450분)다. 941은 대화의 97\%가 37-1011 First-Line Supervisors of Housekeeping and Janitorial Workers(303분)다. 청소 작업 자체가 아니라 청소 감독자 SOC의 대화다. 322는 대화의 99\%가 41-3031 Securities, Commodities, and Financial Services Sales Agents(270분)로 연결된다(연계표 배분 결과이며, 이 연결이 적절한지는 확인하지 않았다). 851은 연계 SOC가 51-4041 Machinists 하나라($d = 0$) $\tau_j$가 그 SOC 값(249분) 그대로다. 851의 $\tau_j$는 $\bar\tau$보다 짧지만, A 순위가 비슷한 직업(95$\sim$100위, $\tau$ 77$\sim$198분)보다 길어서 순위가 14계단 오른다.
\end{itemize}
이 직업들은 모두 A 순위 40$\sim$110위의 대화 점유율이 낮은 직업이다. 순위 격차는 A 순위가 비슷한 직업과의 상대 비교다. 그래서 순위가 촘촘한 이 구간에서는 이웃 직업보다 $\tau_j$가 조금만 길어도 12$\sim$21계단 오른다. A 상위 30에 드는 직업 중 RG가 가장 큰 것은 121 행정 및 경영 지원 관리자(+8)다(표~\ref{tab:top30A}).

\paragraph{A 순위가 높은 직업.} 하위 15개의 $\tau_j$는 24$\sim$135분으로 모두 소분류 중앙값(197분)보다 짧고, 대부분 전체 평균(259분)의 절반 이하다. 다음 경우가 겹친다.
\begin{itemize}[leftmargin=*]
  \item \textbf{판매$\cdot$대면 서비스$\cdot$간병}(522 상품 대여, 521 매장 판매, 532 방문 및 노점 판매, 953 판매 관련 단순, 531 통신 관련 판매, 992 계기$\cdot$자판기 및 주차 관리, 422 요양보호사 및 간병인). 41-2021 Counter and Rental Clerks(90분), 41-2011 Cashiers(58분), 41-9091 Door-to-Door Sales Workers(36분), 31-1131 Nursing Assistants(37분)처럼 대화당 절감 시간이 짧은 SOC로 연결된다. 522와 521은 대화 점유율 상위 30에 드는 직업이지만 B로는 48위, 55위다. 같은 돌봄 직업이라도 보육$\cdot$아동 돌봄(표~\ref{tab:pos})은 B 순위가 높고 간병(422)은 낮다.
  \item \textbf{한 SOC가 대화를 차지하는 경우}. 411 경찰$\cdot$소방 및 교도는 대화의 86\%가 33-3012 Correctional Officers and Jailers(33분)다. 같은 소분류의 33-1012 First-Line Supervisors of Police and Detectives(409분)는 대화의 5.6\%이지만 절감 시간의 41\%를 차지한다. 999 기타 서비스 단순 종사자는 78\%가 39-3093 Locker Room Attendants(24분), 241 의사$\cdot$한의사 및 치과의사는 74\%가 29-1212 Cardiologists(48분)다. 이 직업들의 $\tau_j$는 사실상 그 SOC 하나의 값이다.
  \item \textbf{정보 통신 관련 관리자}(135). 관리자 중 유일하게 하위 15에 들어간다. 연계 SOC가 11-3021 Computer and Information Systems Managers 하나이고, 이 SOC의 $\tau_s$가 115분으로 다른 관리자 SOC(400$\sim$540분)보다 훨씬 짧다.
  \item \textbf{사무}(311 기획 및 마케팅 사무원, 331 회계 및 경리 사무원). 311은 43-4151 Order Clerks(53분), 331은 43-3031 Bookkeeping Clerks(125분)로 연결된다. paper 29에서 ``노출도 > 사용''으로 나온 311과 331은 B로 보면 점유율이 더 낮아진다.
\end{itemize}

\subsection{대화 점유율 상위 30 직업의 B 순위}

\begin{small}
\begin{longtable}{R{0.6cm}L{3.8cm}R{0.9cm}R{0.9cm}R{0.7cm}R{0.7cm}R{0.9cm}R{0.7cm}L{5.8cm}}
\caption{A(대화 점유율) 상위 30 소분류의 B 순위와 대화 1위 SOC(소분류 대화 중 몫, $\tau_s$)}\label{tab:top30A}\\
\toprule
$r$(A) & 소분류 & A(\%) & B(\%) & $r$(B) & RG & DS(\%p) & $\tau$(분) & 대화 1위 SOC \\
\midrule
\endfirsthead
\toprule
$r$(A) & 소분류 & A(\%) & B(\%) & $r$(B) & RG & DS(\%p) & $\tau$(분) & 대화 1위 SOC \\
\midrule
\endhead
% 04_analysis_kor_conv_saved_gap.do가 생성. 직접 고치지 말 것.
1 & 291 작가 및 언론 관련 전문가 & 11.92 & 10.47 & 2 & -1 & -1.45 & 228 & 27-3041 Editors (30.5\%, 241분) \\
2 & 222 컴퓨터 시스템 및 소프트웨어 전문가 & 9.95 & 12.21 & 1 & +1 & +2.26 & 318 & 15-1299 Computer Occupations, All Other (24.2\%, 168분) \\
3 & 225 정보 시스템 및 웹 운영자 & 8.31 & 5.34 & 3 & 0 & -2.98 & 166 & 15-1299 Computer Occupations, All Other (57.8\%, 168분) \\
4 & 372 비서 및 사무 지원 종사자 & 4.69 & 5.24 & 4 & 0 & +0.55 & 290 & 43-9022 Word Processors and Typists (32.6\%, 258분) \\
5 & 283 상품 기획‧홍보 및 조사 전문가 & 4.20 & 5.15 & 5 & 0 & +0.95 & 318 & 27-3043 Writers and Authors (57.3\%, 290분) \\
6 & 223 네트워크 및 정보 보안 전문가 & 3.84 & 3.74 & 7 & -1 & -0.09 & 253 & 15-1299 Computer Occupations, All Other (62.6\%, 168분) \\
7 & 292 학예사‧사서 및 기록물 관리사 & 3.81 & 2.29 & 12 & -5 & -1.52 & 156 & 25-4022 Librarians and Media Collections Specialists (76.3\%, 122분) \\
8 & 213 생명 및 자연과학 관련 시험원 & 2.79 & 3.62 & 8 & 0 & +0.83 & 337 & 15-2051 Data Scientists (34.5\%, 371분) \\
9 & 332 통계‧데이터 관련 사무원 & 2.61 & 4.03 & 6 & +3 & +1.42 & 400 & 15-2051 Data Scientists (36.8\%, 371분) \\
10 & 340 금융 사무 종사자 & 2.61 & 1.78 & 16 & -6 & -0.83 & 177 & 13-2071 Credit Counselors (79.0\%, 192분) \\
11 & 293 연극‧영화 및 영상 전문가 & 2.43 & 2.85 & 10 & +1 & +0.42 & 304 & 27-3023 News Analysts, Reporters, and Journalists (56.4\%, 253분) \\
12 & 211 생명 및 자연과학 관련 전문가 & 2.04 & 2.89 & 9 & +3 & +0.84 & 366 & 19-1029 Biological Scientists, All Other (31.2\%, 459분) \\
13 & 282 금융 및 보험 전문가 & 1.98 & 2.48 & 11 & +2 & +0.50 & 325 & 13-2052 Personal Financial Advisors (20.2\%, 261분) \\
14 & 264 문리‧기술 및 예능 강사 & 1.85 & 2.07 & 13 & +1 & +0.21 & 289 & 25-3041 Tutors (39.8\%, 397분) \\
15 & 235 기계‧로봇공학 기술자 및 시험원 & 1.71 & 1.70 & 17 & -2 & -0.01 & 258 & 17-2141 Mechanical Engineers (43.7\%, 231분) \\
16 & 269 기타 교육 전문가 & 1.41 & 2.02 & 14 & +2 & +0.61 & 372 & 25-9031 Instructional Coordinators (50.8\%, 452분) \\
17 & 131 연구‧교육 및 법률 관련 관리자 & 1.39 & 1.35 & 19 & -2 & -0.04 & 253 & 11-9032 Education Administrators, Kindergarten through Secondary (44.0\%, 232분) \\
18 & 371 총무 사무원 및 대학 행정조교 & 1.33 & 1.01 & 26 & -8 & -0.32 & 197 & 43-9061 Office Clerks, General (75.5\%, 173분) \\
19 & 212 인문 및 사회과학 전문가 & 1.26 & 1.12 & 23 & -4 & -0.14 & 230 & 19-3011 Economists (28.7\%, 291분) \\
20 & 295 디자이너 & 1.19 & 1.01 & 25 & -5 & -0.18 & 220 & 27-1024 Graphic Designers (76.3\%, 182분) \\
21 & 224 데이터 전문가 & 1.10 & 1.82 & 15 & +6 & +0.72 & 430 & 15-1243 Database Architects (87.6\%, 448분) \\
22 & 522 상품 대여 종사자 & 1.07 & 0.37 & 48 & -26 & -0.70 & 90 & 41-2021 Counter and Rental Clerks (100.0\%, 90분) \\
23 & 239 기타 공학 전문가 및 관련 종사자 & 1.07 & 1.33 & 20 & +3 & +0.26 & 323 & 17-2199 Engineers, All Other (35.1\%, 370분) \\
24 & 271 법률 전문가 & 1.04 & 1.25 & 22 & +2 & +0.22 & 313 & 23-1011 Lawyers (71.8\%, 292분) \\
25 & 122 마케팅 및 광고‧홍보 관리자 & 1.01 & 1.50 & 18 & +7 & +0.49 & 385 & 11-2011 Advertising and Promotions Managers (72.2\%, 360분) \\
26 & 261 대학교수 및 강사 & 1.01 & 1.12 & 24 & +2 & +0.11 & 287 & 25-9044 Teaching Assistants, Postsecondary (24.9\%, 297분) \\
27 & 284 감정‧기술 영업 및 중개 관련 종사자 & 1.01 & 0.75 & 31 & -4 & -0.26 & 193 & 41-9031 Sales Engineers (32.7\%, 231분) \\
28 & 510 영업 종사자 & 0.92 & 0.50 & 39 & -11 & -0.42 & 141 & 41-9011 Demonstrators and Product Promoters (43.4\%, 190분) \\
29 & 121 행정 및 경영 지원 관리자 & 0.87 & 1.32 & 21 & +8 & +0.45 & 395 & 11-3131 Training and Development Managers (30.6\%, 468분) \\
30 & 521 매장 판매 종사자 & 0.83 & 0.26 & 55 & -25 & -0.57 & 80 & 41-2011 Cashiers (82.3\%, 58분) \\

\bottomrule
\end{longtable}
\end{small}

대화 점유율 상위 30개 중 B 상위 30에도 드는 직업은 26개다. 상위 10개 중 8개는 B에서도 상위 10에 남는다. 빠지는 둘은 292 학예사$\cdot$사서 및 기록물 관리사(A 7위 $\to$ B 12위, $\tau$ 156분)와 340 금융 사무 종사자(10위 $\to$ 16위, 177분)이고, 211 생명 및 자연과학 관련 전문가(B 9위)와 293 연극$\cdot$영화 및 영상 전문가(B 10위)가 들어온다. 1$\cdot$2위는 서로 바뀐다. A로는 291 작가 및 언론(11.9\%)이 1위지만, B로는 222 컴퓨터 시스템 및 소프트웨어(12.2\%, $\tau$ 318분)가 1위다. 225 정보 시스템 및 웹 운영자는 $\tau$가 166분으로 짧아 점유율이 8.3\%에서 5.3\%로 129개 중 가장 많이 줄지만(DS $-2.98$\%p) 순위는 3위 그대로다.

상위 30 안에서 순위가 크게 내려가는 직업은 522 상품 대여($-26$), 521 매장 판매($-25$), 510 영업 종사자($-11$), 371 총무 사무원 및 대학 행정조교($-8$)다. 크게 올라가는 직업은 121 행정 및 경영 지원 관리자(+8), 122 마케팅 및 광고$\cdot$홍보 관리자(+7), 224 데이터 전문가(+6)다.

\subsection{대분류}

\begin{longtable}{L{4.6cm}R{0.9cm}R{1.3cm}R{1.3cm}R{1.6cm}R{1.4cm}R{1.6cm}R{1.6cm}}
\caption{대분류별 A, B 점유율 합, 대화당 절감 시간, 순위 격차(RG는 소분류 단순 평균)}\label{tab:ksco1}\\
\toprule
대분류 & 소분류 & A(\%) & B(\%) & DS(\%p) & $\tau$(분) & 평균 RG & RG > 0 \\
\midrule
\endfirsthead
% 04_14_analysis_kor_conv_saved_gap.do가 생성. 직접 고치지 말 것.
1 관리자 & 16 & 6.9 & 8.9 & +2.03 & 336 & +8.1 & 13 \\
2 전문가 및 관련 종사자 & 45 & 70.5 & 72.2 & +1.68 & 266 & +1.6 & 28 \\
3 사무 종사자 & 13 & 14.6 & 14.4 & -0.24 & 255 & -3.8 & 3 \\
4 서비스 종사자 & 13 & 1.1 & 0.5 & -0.58 & 122 & -4.2 & 4 \\
5 판매 종사자 & 5 & 3.2 & 1.3 & -1.91 & 103 & -19.2 & 0 \\
6 농림∙어업 숙련 종사자 & 4 & 0.0 & 0.0 & -0.02 & 37 & -3.5 & 0 \\
7 기능원 및 관련 기능 종사자 & 14 & 1.9 & 1.4 & -0.51 & 191 & +0.7 & 7 \\
8 장치∙기계 조작 및 조립 종사자 & 10 & 0.9 & 1.0 & +0.09 & 286 & +4.0 & 8 \\
9 단순노무 종사자 & 9 & 0.9 & 0.4 & -0.52 & 105 & -4.4 & 4 \\

\bottomrule
\end{longtable}

대분류 $\tau$는 대분류 안 $\sum S / \sum C$다. 관리자(336분), 장치$\cdot$기계 조작(282분), 전문가(266분), 사무(255분)는 전체 평균(259분) 근처 이상이고, 판매(103분), 단순노무(105분), 서비스(122분), 농림어업(33분)은 절반 이하다. A에서 B로 바꾸면 관리자와 전문가의 몫이 합쳐 3.7\%p 늘고, 판매$\cdot$서비스$\cdot$단순노무의 몫이 합쳐 3.0\%p 준다. 장치$\cdot$기계 조작은 10개 중 9개가 RG > 0이지만 점유율이 작아(1\% 안팎) 점유율 차이는 0.08\%p에 그친다.

\section{해석과 유의사항}

\begin{enumerate}[leftmargin=*]
  \item \textbf{A와 B 중 무엇을 쓸지는 상위 직업에는 거의 영향이 없다.} 상위 30 중 26개가 겹치고 순위상관이 0.96이다. 차이는 판매$\cdot$대면 서비스 직업에서 크다. 일자리 영향 시나리오(paper 27의 점유율 $\times$ X명)에서 이 직업들은 B$\cdot$C를 쓰면 A의 절반 이하 몫을 받는다. 대화 수는 많아도 대화가 짧은 과업이라는 AEI의 추정 때문이다.
  \item \textbf{$\tau_j$의 차이는 한국 직업의 차이가 아니라 SOC 구성의 차이다.} $\tau_s$는 전 세계 Claude.ai 대화의 평균이다. 같은 SOC라면 한국 대화도 같은 시간을 줄인다고 가정했다. 한국 대화의 과업 구성이 전 세계와 다르면 $\tau_j$도 달라질 수 있다.
  \item \textbf{SOC 하나가 대화를 차지하는 소분류는 연계 경로에 민감하다.} 411(Correctional Officers 86\%), 999(Locker Room Attendants 78\%), 241(Cardiologists 74\%), 522$\cdot$135(SOC 하나)는 $\tau_j$가 SOC 하나의 값이다. 241에서 대화의 대부분이 의사 SOC 중 Cardiologists 하나로 연결되는 이유는 이 노트에서 확인하지 않았다(AEI의 SOC는 사용자의 직업이 아니라 대화 과업의 직업이므로, AEI 과업 분류와 연계표 배분을 함께 확인해야 한다). 이 직업들의 순위 격차는 연계표 배분(paper 21)과 AEI 과업 분류에 따라 달라질 수 있다.
  \item \textbf{절감 시간은 추정치다.} $\tau_s$는 Claude가 ``사람이 혼자 하면 걸릴 시간''과 ``AI와 함께 걸린 시간''을 대화마다 추정한 값의 차이다\parencite{tamkin-2025-estimating-ai-productivity-gains-from}. 실측이 아니다. 따라서 관리자$\cdot$교육 직업의 높은 $\tau_s$와 판매$\cdot$서비스 직업의 낮은 $\tau_s$에는 모델 추정의 오차가 섞여 있을 수 있다.
  \item \textbf{순위 격차는 분포의 가운데에서 커진다.} 순위가 촘촘한 A 44$\sim$86위 구간의 $|\mathrm{RG}|$ 평균(11.7)은 위$\cdot$아래 구간(각 5.3)의 두 배다. 상위 15 목록에는 대화 점유율이 낮은 직업이 많이 들어온다. 강도 자체를 비교하려면 표의 $\tau_j$(= LR 순위)를 함께 본다.
\end{enumerate}

\section{산출 파일}

\begin{itemize}[leftmargin=*]
  \item 코드: \texttt{code/04\_14\_analysis\_kor\_conv\_saved\_gap.do}(Stata 17). 입력은 \texttt{kor\_ai\_usage\_ksco3}(paper 25$\cdot$27), \texttt{kor\_soc6\_conv\_ai\_indicators}, SOC 2018--KSCO 8차 연계표와 SOC 명칭. 기여 합이 소분류 A, B와 같음을 assert로 확인한다.
  \item 표: \texttt{results/table/31\_kor\_conv\_saved\_gap.xlsx}. 시트 \texttt{ksco3}(129개 소분류의 A, B, 순위, RG, DS, LR, $\tau$, 노출도, 대화 1위$\cdot\tau$ 최대 상승$\cdot$최대 하락 SOC), \texttt{top}(RG 상위$\cdot$하위 15), \texttt{top30A}(A 상위 30), \texttt{soc}(상위$\cdot$하위 15 소분류 $\times$ SOC의 대화 몫, 절감 시간 몫, $d_{sj}$, $\tau_s$), \texttt{ksco1}, \texttt{stats}.
  \item 이 노트의 표: \texttt{results/table/31\_summary.tex}, \texttt{31\_top\_pos.tex}, \texttt{31\_top\_neg.tex}, \texttt{31\_top30A.tex}, \texttt{31\_ksco1.tex}.
\end{itemize}

\printbibliography[title={참고문헌}]

\end{document}
